% paper/sections/design.tex — Software Design and Architecture
% Written in Phase 88 (Plan 04). TikZ architecture diagram added in Phase 90.

\section{Software Design and Architecture}
\label{sec:design}

\subsection{Layered Architecture}

\begin{figure}[ht]
\centering
\begin{tikzpicture}[
  font=\small,
  layer/.style={draw, rounded corners=4pt, minimum width=10.5cm, minimum height=1.35cm,
                align=center, inner sep=6pt},
  arr/.style={{Stealth[length=5pt]}-{Stealth[length=5pt]}, thick},
  crosslbl/.style={font=\footnotesize, align=left, fill=white, inner sep=2pt},
]
  % Layer boxes (stacked top-to-bottom); descriptions live INSIDE each box
  % so they never collide with the boundary-crossing arrow labels.
  \node[layer, fill=blue!10] (python) at (0,0)
    {\textbf{Python API layer} (\texttt{fdars})\\[2pt]
     {\footnotesize\itshape Fdata container $\cdot$ sklearn estimators $\cdot$ grounded advisor $\cdot$ dataset loaders}};
  \node[layer, fill=orange!10, below=2cm of python] (pyo3)
    {\textbf{PyO3 binding layer} (\texttt{fdars.\_native})\\[2pt]
     {\footnotesize\itshape \texttt{\#[pyfunction]} wrappers $\cdot$ \texttt{PyReadonlyArray} views $\cdot$ row\,$\leftrightarrow$\,col-major}};
  \node[layer, fill=green!10, below=2cm of pyo3] (core)
    {\textbf{Rust core} (\texttt{fdars-core})\\[2pt]
     {\footnotesize\itshape \texttt{FdMatrix} (nalgebra) $\cdot$ data-parallel loops (rayon) $\cdot$ algorithm modules}};

  % One bidirectional arrow per boundary, with a single centred label.
  \draw[arr] (python) -- (pyo3)
    node[crosslbl, midway] {NumPy arrays $\leftrightarrow$ Python objects};
  \draw[arr] (pyo3) -- (core)
    node[crosslbl, midway] {\texttt{FdMatrix} $\leftrightarrow$ results (\texttt{convert.rs})};
\end{tikzpicture}
\caption{Three-layer architecture of \texttt{fdars}.  \texttt{PyReadonlyArray}
wrappers give the Rust side a read-only view of the caller's NumPy arrays; 2D
arrays are then copied once into the column-major layout that \texttt{FdMatrix}
expects (\texttt{src/convert.rs}).}
\label{fig:architecture}
\end{figure}

\texttt{fdars} is organised as three cooperating layers
(Figure~\ref{fig:architecture}).

\textbf{Computation core.}  Numerical algorithms live in the external Rust
crate \texttt{fdars-core}, which provides the \texttt{FdMatrix} linear-algebra
abstraction backed by \texttt{nalgebra} and data-parallel loops via
\texttt{rayon}.  \texttt{fdars-core} 0.33.0 is the minimum required version
declared in \texttt{Cargo.toml}, compiled with the \texttt{parallel} feature
to enable rayon-backed parallelism.  No Python objects cross into this layer;
all communication occurs through the binding layer described next, and
\texttt{fdars-core} exports no Python-facing symbols.

\textbf{PyO3 binding layer.}  A set of 27 Rust modules (one per
functional category) exposes \texttt{fdars-core} algorithms to Python via
\href{https://pyo3.rs}{PyO3} 0.28.  Each binding function is annotated
\texttt{\#[pyfunction]} and accepts NumPy arrays through
\texttt{PyReadonlyArray} wrappers, which hold a read-only view of the
caller's array for the duration of the Rust call without an extra data copy;
because the bindings keep the GIL held for the duration of each call, Python
code cannot mutate the array while Rust reads it.
Default parameter values are declared with \texttt{\#[pyo3(signature = (...))]}
at the function level, allowing idiomatic Python keyword-argument syntax at
call sites.  All modules are registered via a shared macro
(\texttt{register\_submodule!}) in \texttt{src/lib.rs} that creates a
\texttt{PyModule}, calls the module's \texttt{register} function, and attaches
the submodule to the parent extension --- so adding a new algorithm family
requires a new Rust module (declared and registered in \texttt{src/lib.rs})
and an entry in the package initialiser's submodule list.  Together with three
pure-Python modules (\texttt{covariance}, \texttt{datasets},
\texttt{metrics}), the package exposes \nsubmodules{} public analysis
submodules (in addition to the \texttt{advisor}, \texttt{sklearn},
\texttt{plot}, \texttt{mcp}, and \texttt{results} packages).

\textbf{Boundary convention.}  The PyO3 boundary involves a layout conversion:
NumPy stores two-dimensional arrays in C (row-major) order, while
\texttt{fdars-core}'s \texttt{FdMatrix} is column-major.  The conversion
function \texttt{numpy2d\_to\_fdmatrix} (see \texttt{src/convert.rs}) performs
an explicit element-wise transpose: given an $(n_\mathrm{obs} \times
n_\mathrm{points})$ NumPy array in C order it builds the column-major flat
vector required by \texttt{FdMatrix::from\_column\_major}.  The reverse
function \texttt{fdmatrix\_to\_numpy2d} calls \texttt{to\_row\_major()} before
allocating the output array.  The \texttt{PyReadonlyArray} wrapper avoids an
extra data copy for the view itself, but a layout conversion is always present
at the 2D matrix boundary.  Errors from \texttt{fdars-core} propagate as Rust
\texttt{Result<T, FdarError>} values and are converted to Python
\texttt{ValueError} exceptions by the helper \texttt{to\_pyresult} in
\texttt{src/convert.rs}, so users receive descriptive error messages rather
than opaque panics.

\textbf{Python API layer.}  On top of the native extension, a pure-Python
layer provides idiomatic Python-level conveniences: the \texttt{Fdata}
object-oriented container (Section~\ref{sec:represent}), the
\texttt{cluster\_optim} parameter-search orchestrator, the grounded advisor
module (Section~\ref{sec:advisor}), plotting utilities (\texttt{fdars.plot}),
ergonomic dataset loaders (\texttt{fdars.datasets}), and the scikit-learn
estimator wrappers described in Section~\ref{subsec:sklearn}.  The module
package initialiser registers every native submodule and adds pure-Python
helpers to its namespace, so both \texttt{fdars.depth.fraiman\_muniz\_1d} and
\texttt{from fdars.depth import fraiman\_muniz\_1d} work as users expect.

\subsection{Module Map}

\texttt{fdars} exposes \nsubmodules{} public analysis submodules (in
addition to the \texttt{advisor}, \texttt{sklearn}, \texttt{plot},
\texttt{mcp}, and \texttt{results} packages); the capability map
bundled with the package catalogues \npubliccallables{} public module
callables, organised by method family below (the pure-Python
\texttt{datasets}, \texttt{covariance} and \texttt{metrics} helper modules
provide dataset loaders, covariance kernels and evaluation metrics):

\begin{description}
  \item[Representation, basis and smoothing]
    \texttt{fdata}, \texttt{represent}, \texttt{basis}, \texttt{smoothing} ---
    functional object construction, derivative estimation, basis expansion,
    GCV-penalised smoothing, interpolation, resampling;
    \texttt{multi\_fdata} --- an opaque container (\texttt{PyMultiFunData})
    that groups several functional components, each with its own evaluation
    grid, over the same set of observations; at present it can only be
    inspected (\texttt{n\_obs}, \texttt{n\_components}), and no analysis
    routine accepts it as input.

  \item[Alignment]
    \texttt{alignment} --- elastic registration via the square-root slope
    function (SRSF) framework \citep{srivastava_et_al_2011}, including
    pairwise warping and Karcher mean under the elastic metric.

  \item[Depth, outliers and scoring]
    \texttt{depth}, \texttt{outliers} --- Fraiman--Muniz
    depth \citep{fraiman_muniz_2001}, band depth
    \citep{lopez_pintado_romo_2009}, an outliergram-style MEI/MBD shape-outlier
    detector adapted from \citet{arribas_gil_romo_2014} (the parabola is fitted
    to the data by least squares rather than taken from the published
    theoretical bound), a Fast-MUOD-style detector adapted from
    \citet{ojo_fast_muod_2022} (itself a variant of MUOD,
    \citealp{azcorra_muod_2018}) that uses the pointwise mean rather than the
    pointwise median as the reference curve, and a magnitude/shape outlyingness
    decomposition (magnitude $= 1 -$ modified band depth, shape $=$ distance
    of the normalised centred curve from the mean direction) that returns
    scores only and is in the spirit of, but not identical to, the MS-plot of
    \citet{dai_genton_2018};
    \texttt{scoring} --- functional prediction-error metrics (MSE, MAE, MAPE,
    MSLE, explained variance).

  \item[FPCA and sparse/longitudinal PACE]
    \texttt{regression.fpca} (also available as \texttt{Fdata.to\_pc}),
    \texttt{pace\_fpca} ---
    dense FPCA and the PACE algorithm for irregularly observed trajectories
    \citep{yao_muller_wang_2005}.

  \item[Clustering and classification]
    \texttt{clustering}, \texttt{classification}, \texttt{shapelet} ---
    functional $k$-means \citep{tarpey_kinateder_2003}, fuzzy $c$-means,
    Gaussian mixtures, elastic-alignment clustering, DBSCAN, and FunFEM \citep{bouveyron_et_al_2015};
    FPCA-based $k$-NN, LDA and QDA, kernel and DD-classifiers, an elastic
    multinomial classifier, and shapelet discovery and transform
    \citep{ye_keogh_2009}.

  \item[Regression]
    \texttt{regression}, \texttt{scalar\_on\_function}, \texttt{famm} ---
    scalar-on-function (FPC-based and non-parametric), function-on-function
    regression, and functional additive mixed models.

  \item[Functional time series and seasonal decomposition]
    \texttt{fts} --- functional time-series modelling via the FTSM framework
    \citep{hyndman_ullah_2007}: FPCA of the curve sequence followed by an
    independent univariate autoregression on each score series, with
    multi-step forecasting;
    \texttt{seasonal} --- period estimation, peak detection, seasonal-strength
    and seasonality-change measures, and STL/SSA decomposition of functional
    data.

  \item[Statistical process monitoring]
    \texttt{spm} --- phase-1 baseline estimation and phase-2 monitoring via
    functional $T^2$ and SPE statistics.

  \item[Conformal prediction and tolerance bands]
    \texttt{conformal}, \texttt{tolerance} --- FPCA-based tolerance bands and
    conformal scalar-on-function prediction intervals.

  \item[Density, Fr\'{e}chet and metric]
    \texttt{density\_fda}, \texttt{frechet}, \texttt{metric} ---
    log-quantile-density FPCA \citep{petersen_muller_2016}, Wasserstein
    barycenter \citep{agueh_carlier_2011}, Fr\'{e}chet means and regression
    \citep{petersen_muller_2019}, and $L^p$ and DTW \citep{sakoe_chiba_1978} distance matrices.

  \item[Inference, explanation and simulation]
    \texttt{inference}, \texttt{explain}, \texttt{simulation} ---
    functional ANOVA, permutation and bootstrap tests; model explanation
    (FPC-level SHAP, LIME, ALE/PDP, permutation importance, counterfactuals,
    Sobol indices, influence diagnostics); Gaussian-process and
    Karhunen--Lo\`{e}ve (basis-expansion) simulation.
\end{description}

This organisation mirrors the taxonomy used in the comparison table
(Table~\ref{tab:comparison}); see Section~\ref{sec:comparison} for a
side-by-side evaluation of method coverage across peer packages.

\subsection{The \texttt{Fdata} Container}
\label{subsec:fdata}

At the Python level, dense functional data are held in the \texttt{Fdata}
container, which bundles the observation matrix with its evaluation grid,
identifiers, and sample-level metadata and exposes \nfdata{} public methods,
most of which delegate to the binding layer.  Section~\ref{sec:represent} describes it in
full.

\subsection{Scikit-learn Estimator Layer}
\label{subsec:sklearn}

A pure-Python scikit-learn-compatible layer (\texttt{fdars.sklearn}, installed
with the \texttt{[sklearn]} extra) wraps the most commonly used bindings as
estimator classes.  \nsklearnestimators{}
estimators pass the full \texttt{check\_estimator} battery
\citep{pedregosa_sklearn_2011}, covering smoothing and basis transformers,
FPCA-based regression and classification, depth-based classifiers and outlier
detectors, and functional clustering estimators (e.g., \texttt{FPCATransformer},
\texttt{FPCRegressor}, \texttt{FunctionalKMeans}, \texttt{LRTOutlierDetector}).

Conformance with \texttt{check\_estimator} carries concrete guarantees: the
estimator's \texttt{get\_params}/\texttt{set\_params} round-trip correctly
(enabling hyperparameter grid search), \texttt{clone} produces an independent
copy, and each estimator is picklable (required for multi-process cross-validation
via \texttt{GridSearchCV(n\_jobs=-1)}).  These estimators therefore integrate
without modification into \texttt{sklearn.pipeline.Pipeline} and
cross-validation utilities, as illustrated in the capability tour
(Section~\ref{sec:capabilities}).

Some methods in \texttt{fdars} are deliberately excluded from the estimator
layer (\nexcludedmethods{} methods, listed in \texttt{fdars.sklearn.EXCLUDED\_METHODS})
because they cannot satisfy the scikit-learn estimator contract --- for
example, paired elastic alignment
(\texttt{alignment.elastic\_align\_pair}) and the Karcher mean
(\texttt{alignment.karcher\_mean}) produce batch-dependent outputs that fail
scikit-learn's subset-invariance check, and the PACE FPCA
(\texttt{pace\_fpca.pace\_fpca}) requires ragged per-observation grids
rather than a rectangular matrix.  This boundary is documented transparently
so that users building custom pipelines know which methods to call directly.
